%%%%%%%%%%%%%%%%%%%%%%%%%%%%%%%%%%%%%%%%%%%%%%%%%%%%%%%%%%%%%%%%%%%%%%%%%%%
% CHEP 2026 proceedings -- CernVM-FS Filebundles
% EPJ Web of Conferences (Web of Conferences LaTeX template)
%%%%%%%%%%%%%%%%%%%%%%%%%%%%%%%%%%%%%%%%%%%%%%%%%%%%%%%%%%%%%%%%%%%%%%%%%%%

\documentclass{webofc}
\usepackage[varg]{txfonts}   % Web of Conferences font
\usepackage{hyperref}
\usepackage{url}
\usepackage{listings}
\hypersetup{colorlinks=true,citecolor=blue,urlcolor=blue,linkcolor=blue}

\newcommand{\cvmfs}{CernVM-FS}
\newcommand{\fuse}{FUSE}
\newcommand{\wlcg}{WLCG}
% Compact listing style for the JSON bundle example
\lstdefinestyle{bundle}{
  basicstyle=\ttfamily\footnotesize,
  breaklines=true,
  columns=fullflexible,
  frame=single,
  framesep=4pt,
  xleftmargin=4pt,
  showstringspaces=false,
}

\begin{document}

\title{CernVM-FS Filebundles for low-latency software distribution
in interactive use cases}

\author{%
  \firstname{Valentin} \lastname{Volkl}\inst{1}\fnsep\thanks{\email{valentin.volkl@cern.ch}} \and
  \firstname{Georgios} \lastname{Christodoulis}\inst{1} \and
  \firstname{Jakob} \lastname{Blomer}\inst{1}
}

\institute{CERN, 1211 Geneva 23, Switzerland}

\abstract{%
\cvmfs\ is a read-only, on-demand file system used to distribute
scientific software to the Worldwide LHC Computing Grid and many
communities beyond high-energy physics. Its lazy, per-file fetching is
highly efficient for batch computing, but it adds noticeable latency to
interactive sessions: starting an application that opens hundreds of
small files on a cold cache serializes one network round trip after
another. We present \emph{filebundles}, a mechanism that exploits the
predictable file-access pattern at application startup. A small bundle
specification, published next to a trigger file, declares the set of
files needed immediately after the trigger is opened. When a client
opens the trigger, \cvmfs\ prefetches the listed dependencies
concurrently using a pool of worker threads, replacing a long serial
chain of requests with a short parallel prefetch phase. On a ROOT
startup benchmark, filebundles bring performance into the local-disk
range and keep startup time nearly flat as network latency grows, while
the on-demand baseline degrades approximately linearly. Filebundles are
an opt-in feature in \cvmfs\ 2.14.
}

\maketitle

%======================================================================
\section{Introduction}
\label{sec-intro}

\cvmfs~\cite{ref-cvmfs} is a global, read-only file system for
distributing scientific software. 
It is a core component of the Worldwide LHC Computing Grid ( \wlcg\ )and
is also used by communities beyond high-energy physics, including
Euclid, LIGO, LSST, SKA, and EESSI.

\cvmfs\ uses a distributed architecture (Fig.~\ref{fig-hierarchy}), where software is originally published in Stratum0 - a server which is considered the source of truth, in the form of repositories.
The contents of Stratum0 are then cloned in a series of geographically scattered mirror servers (Stratum1's), arranged to facilitate the different user sites, in order to mitigate the traffic of Stratum0.
Users can access the contents of a repository either directly from a Stratum (0 or 1) (not encouraged) or via a series of forward proxies that site administrators configure, for better load balancing and caching.
\wlcg\ has about fifteen Stratum~1 mirror servers in Europe, the Americas, and Asia, managing about $3.3\times10^{10}$ files
and 2~PB of data and more than 8000 container images.

Data (repository contents) on \cvmfs\ are stored in a content addressable format.
During publication files are chunked, hashed and compressed.
A filename is associated to its hash (or hashes in the case of chunking) through dedicated dictionaries called catalogs.
Catalogs can also be nested in order to sustain the rapid repository growth.

From the users perspective, repositories can be accessed via a global namespace, by default under /cvmfs in their machine.
This is done by the \cvmfs\ client, which is a \fuse\ based filesystem.
Typically, there is some disk allocated space in the user's machine that works as a local cache.
\cvmfs\ is performing cache cleanup using an LRU policy.

For simple listing operations, the client will fetch the associated catalogs which contain all the information necessary.
When an \emph{open()}  is issued from the user, the client will lookup for the correct hash, and will issue a \emph{fetch()}.
During fetching, the client will download the compressed hash, will decompress it and store it in the local cache, available for reading.
On a \emph{read()} the client will make the contents of the files available to the user immediately, or will trigger more \emph{fetch()} operations in case the file is chunked, or the hash has already been evicted from the cache.
According to this mechanism, on first access the cache is empty (cold cache) and the contents of the repository become available after they are needed.
Downloading, decompression, and caching are transparent to the user.


This lazy, content-addressed model is extremely effective for
distributed batch computing, where the benefits of a shared software
distribution system dominate any per-file overhead. In interactive
sessions, however, the same on-demand behaviour can make application
startup noticeably slower than from local disk. This paper presents
\emph{filebundles}, a mechanism that reduces that startup latency by
prefetching the files an application needs immediately after launch.

\begin{figure}[t]
\centering
\includegraphics[width=0.78\textwidth]{figures/cvmfs_hierarchy.png}
\caption{The \cvmfs\ cache hierarchy: a writable Stratum~0 origin,
globally distributed Stratum~1 mirrors, site-level Squid proxies, and
per-client disk and kernel page caches. Filebundles change only the
timing of the first wave of requests, not the data model.}
\label{fig-hierarchy}
\end{figure}
\clearpage % Keep the Section 2 hierarchy figure out of the bundle listing.

%======================================================================
\section{Motivation}
\label{sec-motivation}

At the CernVM Users Workshop in 2022, Chris Burr (LHCb) reported that
interactive \cvmfs\ use can be roughly three to five times slower
than local disk, even with a hot cache~\cite{ref-burr}. This overhead is
usually acceptable for distributed batch workloads, but is noticeable
when starting an interactive analysis environment.

The effect is easy to reproduce on a controlled testbed. Installing
ROOT from a release tarball on local disk, importing ROOT in Python
takes about 2.47~s. Using the same software stack from \texttt{/cvmfs}
with a cold cache, the time increases to about 4.69~s. In this setup
the network path is very short, so the slowdown is not primarily a
bandwidth problem. It confirms that a significant part of the delay
comes from the way \cvmfs\ serves many files on demand.

\subsection{Where does the time go?}
\label{sec-time}

When an application such as ROOT starts, it opens many files one after
another. On a cold cache, for each file \cvmfs\ handles a FUSE
upcall, issues an HTTP request, waits for the round trip, downloads the
object, decompresses it and returns the data. This per-file overhead is
paid again and again, and because the requests are dependent they are
effectively serialized. The total amount of work is unremarkable, but
the serialized per-file latency adds up. In batch computing this
trade-off is acceptable; in interactive use the serialization is
precisely what hurts.

It is useful to recall the \cvmfs\ cache hierarchy
(Fig.~\ref{fig-hierarchy}) in which any optimization has to operate. At
the top is the Stratum~0, the writable origin managed by the
repository maintainer (``librarian''). Below it are the globally
distributed Stratum~1 read-only mirrors. Sites typically add a Squid
forward proxy to absorb repeated requests from many clients. On each
worker node or workstation \cvmfs\ keeps a local disk cache, and
above that the Linux kernel maintains the page cache. Filebundles
operate within this existing hierarchy. They do not change the data
model; they only change \emph{when} and \emph{how} the first wave of
file requests is issued.

%======================================================================
\section{Filebundles}
\label{sec-design}

Filebundles use prior knowledge of application startup. For applications
such as ROOT~\cite{ref-root} or Python, the files needed immediately
after a given \emph{trigger} file is opened are often predictable. A
bundle records these dependencies instead of leaving the client to
discover them one file at a time.

A bundle is declared by a specification file published next to its
trigger file in the same directory, named
\texttt{.cvmfsbundle-<filename>}, where \texttt{<filename>} is the
trigger. The specification is a small, versioned JSON document holding
minimal metadata and a flat list of dependency paths, expressed as
absolute paths from the repository root. A representative ROOT bundle,
keyed on the Python package \texttt{\_\_init\_\_.py}, contains roughly
250 paths:

\begin{lstlisting}[style=bundle]
#%CVMFS_BUNDLE version=1 encoding=UTF-8
{
  "name": "CVMFS_BUNDLE",
  "version": "1.0.0",
  "encoding": "UTF-8",
  "dependencies": [
    "/root/lib/libCore.so",
    "/root/lib/libRIO.so",
    "/root/lib/libcppyy.so",
    "/root/lib/ROOT/_application.py",
    ...  (~250 paths total)
  ]
}
\end{lstlisting}

When a client opens the trigger file, \cvmfs\ detects the matching
\texttt{.cvmfsbundle} specification and begins prefetching the listed
dependencies concurrently, using a pool of worker threads, while the
application continues to run. In effect, a long serial chain of
dependent requests is replaced by a much shorter parallel prefetch
phase that warms the local cache ahead of the application's subsequent
\texttt{open()} calls.

The feature is disabled by default and can be enabled per mount with
\texttt{CVMFS\_PREFETCH\_FILEBUNDLES=yes}~\cite{ref-docs}. The setting
is read at mount time; mounts for which it is disabled do not start
bundle prefetching. The prefetch-pool size is controlled by
\texttt{CVMFS\_BUNDLE\_POOL\_SIZE}.

\subsection{Publishing a filebundle}
\label{sec-publish}

From the librarian's point of view the workflow is straightforward.
First, choose a suitable trigger file, for example a Python
\texttt{\_\_init\_\_.py} or a launcher binary such as
\texttt{root.exe}. Second, determine which files are typically read
after that trigger; this can be done by hand or with tracing tools such
as \texttt{strace} or the built-in \texttt{CVMFS\_TRACEFILE}
mechanism. Third, publish the resulting \texttt{.cvmfsbundle-<filename>}
specification next to the trigger file, just like any other file in the
repository. Removing the specification (and republishing) removes the
bundle, after which the trigger behaves like an ordinary file.

%======================================================================
\section{Results}
\label{sec-results}

Unless noted otherwise, the benchmarks use ROOT~6.36.12 in a test
\cvmfs\ repository served from \texttt{localhost}. The client ran on
an AMD EPYC~7302 system (2$\times$16 cores, 64 threads at
${\sim}3.0$~GHz, 251~GB RAM), with AlmaLinux~10.1 and kernel~6.12.
Using \texttt{localhost} removes external network effects, exposing the
per-file software overhead.

\subsection{Startup performance}
\label{sec-startup}

Figure~\ref{fig-violin} compares ROOT startup times on local disk, on
\cvmfs\ with a cold cache and no filebundles, and on \cvmfs\ with
filebundles enabled. In this benchmark, filebundles bring \cvmfs
startup time to the local-disk range. Some measurements are marginally
faster than local disk, consistent with overlap between downloading and
decompression in the prefetch path.

\begin{figure}[t]
\centering
\includegraphics[width=0.62\textwidth]{figures/violin_bundle.pdf}
\caption{Distribution of ROOT startup times. With filebundles enabled
(\texttt{CVMFS\_PREFETCH\_FILEBUNDLES=ON}), \cvmfs\ startup matches
local-disk performance, closing the gap to the on-demand cold-cache
baseline.}
\label{fig-violin}
\end{figure}

\subsection{Prefetch thread pool}
\label{sec-pool}

Prefetching is implemented with a pool of worker threads whose size is
set by \texttt{CVMFS\_BUNDLE\_POOL\_SIZE}. As shown in
Fig.~\ref{fig-pool}, startup performance improves as workers are added
but saturates around eight threads, which is also the default. Most of
the benefit is therefore obtained with a moderate amount of concurrency
and without aggressive tuning.

\begin{figure}[t]
\centering
\includegraphics[width=0.62\textwidth]{figures/pool_scan.pdf}
\caption{ROOT startup time as a function of the prefetch thread-pool
size (\texttt{CVMFS\_BUNDLE\_POOL\_SIZE}). Performance saturates around
eight worker threads, the default.}
\label{fig-pool}
\end{figure}

\subsection{Network latency}
\label{sec-latency}

We injected loopback round-trip latency with \texttt{tc qdisc netem}.
Without filebundles, startup time grows approximately linearly with
round-trip time because the application waits for many dependent,
small-file fetches. With filebundles, startup time remains nearly flat
over a substantial range (Fig.~\ref{fig-latency}): concurrent prefetch
pays the round-trip penalty for a small number of batches rather than
for each file. This is an idealized loopback experiment.

\begin{figure}[t]
\centering
\includegraphics[width=0.62\textwidth]{figures/latency_scan.pdf}
\caption{ROOT startup time as a function of injected round-trip
latency. The on-demand baseline degrades roughly linearly, while
filebundles keep startup nearly flat by overlapping requests.}
\label{fig-latency}
\end{figure}

\subsection{Caveats}
\label{sec-caveats}

Filebundles do not reduce the total amount of data transferred; they
compress the startup traffic into a short, intense burst. In the
benchmark, requests peak at about 1480~per~second with filebundles,
versus about 680~per~second without them (Fig.~\ref{fig-burst}). The
feature should therefore be used selectively, in particular on
interactive systems such as \texttt{lxplus}, where startup latency
matters most. Over-prefetching should be avoided, and any rollout
should be coordinated with proxy and site administrators so that the
request bursts do not overload shared caches.

\begin{figure}[t]
\centering
\includegraphics[width=0.62\textwidth]{figures/req_burst.pdf}
\caption{Request rate during startup. Filebundles concentrate the same
total traffic into a shorter, more intense burst (peak ${\sim}1480$
req/s versus ${\sim}680$ req/s without bundles).}
\label{fig-burst}
\end{figure}

%======================================================================
\section{Outlook}
\label{sec-outlook}

FUSE passthrough in \cvmfs~2.14 on recent kernels can hand the
kernel a descriptor for the cached backing file at \texttt{open()}, so
subsequent reads bypass FUSE. Its interaction with filebundles still
needs to be benchmarked. Planned work for \cvmfs~2.15 includes tools
for creating and updating bundles, support for file chunks in bundle
specifications, and integration with build tools. Packfiles, which
would combine small backend objects to reduce request counts, are also
under consideration; backward compatibility is the principal obstacle.

%======================================================================
\section{Conclusions}
\label{sec-conclusions}

Filebundles reduce application-startup latency when a known set of files
is needed immediately after a trigger is opened. They require the
repository maintainer to trace startup behaviour and publish a
\texttt{.cvmfsbundle} specification next to the trigger. In the ROOT
benchmark, they reduce cold-start time to the local-disk range and make
startup less sensitive to added round-trip latency. Filebundles are an
opt-in \cvmfs~2.14 feature, enabled per mount with
\texttt{CVMFS\_PREFETCH\_FILEBUNDLES}. They are best suited to
interactive systems, where the prefetch burst can be managed.

\vspace{0.5em}
\noindent\textbf{Acknowledgements.} We thank Chris Burr (LHCb) for the
input that motivated this work, and the \cvmfs\ team at CERN.

\vspace{0.3em}
\noindent\emph{Disclosure: AI-based tools were used to assist with
drafting and formatting this manuscript; all content was reviewed and
verified by the authors.}

%======================================================================
\begin{thebibliography}{9}

\bibitem{ref-cvmfs}
J. Blomer, C. Aguado-S\'anchez, P. Buncic, A. Harutyunyan,
Distributing LHC application software and conditions databases using
the CernVM file system, J.\ Phys.: Conf.\ Ser.\ \textbf{331}, 042003
(2011). \url{https://doi.org/10.1088/1742-6596/331/4/042003}

\bibitem{ref-burr}
C. Burr, \textit{Client performance}, CernVM Users Workshop (2022).
\url{https://indico.cern.ch/event/1079490/contributions/4939521/attachments/2506388/4306645/2022-09-12_CVMFS-Workshop_LHCb.pdf}

\bibitem{ref-root}
R. Brun, F. Rademakers, ROOT---an object oriented data analysis
framework, Nucl.\ Instrum.\ Methods Phys.\ Res.\ A \textbf{389},
81--86 (1997). \url{https://doi.org/10.1016/S0168-9002(97)00048-X}

\bibitem{ref-docs}
\cvmfs\ documentation. \url{https://cvmfs.readthedocs.io/}

\end{thebibliography}

\end{document}
